
\documentclass{article}
\usepackage{colortbl}
\usepackage{makecell}
\usepackage{multirow}
\usepackage{supertabular}

\begin{document}

\newcounter{utterance}

\twocolumn

{ \footnotesize  \setcounter{utterance}{1}
\setlength{\tabcolsep}{0pt}
\begin{supertabular}{c@{$\;$}|p{.15\linewidth}@{}p{.15\linewidth}p{.15\linewidth}p{.15\linewidth}p{.15\linewidth}p{.15\linewidth}}

    \# & $\;$A & \multicolumn{4}{c}{Game Master} & $\;\:$B\\
    \hline 

    \theutterance \stepcounter{utterance}  

    & & \multicolumn{4}{p{0.6\linewidth}}{\cellcolor[rgb]{0.9,0.9,0.9}{%
	\makecell[{{p{\linewidth}}}]{% 
	  \tt {\tiny [A$\langle$GM]}  
	 State giocando a un gioco di negoziazione in cui dovete accordarvi su come dividere una serie di oggetti tra voi e un altro giocatore.\\ \tt \\ \tt Regole:\\ \tt (a) A Lei e all'altro giocatore viene dato un insieme di oggetti. Ognuno di voi riceve anche una funzione di valore segreta, che rappresenta il valore di ciascun tipo di oggetto.\\ \tt (b) Si scambiano messaggi con l'altro giocatore per concordare chi si aggiudica gli oggetti. Potete inviare un massimo di 5 messaggi ciascuno, oppure terminare in anticipo facendo una proposta segreta in qualsiasi momento.\\ \tt (c) A ciascuno di voi viene chiesto di inviare una proposta segreta indicando gli oggetti che desiderate, formattata tra parentesi quadre come segue: "[Proposta: <numero> <nome oggetto>, <numero> <nome oggetto>, <...>]".\\ \tt (d) Se le vostre proposte sono complementari, cioè ci sono abbastanza oggetti per soddisfare entrambe le proposte, a ciascun giocatore viene assegnato un punteggio basato sulla somma dei valori degli oggetti ricevuti. In caso contrario, entrambi ricevono zero punti.\\ \tt (e) L'obiettivo è massimizzare il proprio punteggio.\\ \tt \\ \tt Cominciamo.\\ \tt \\ \tt L'insieme degli oggetti disponibili è:\\ \tt \\ \tt 4 carte, 2 orologi, 2 patate\\ \tt \\ \tt La funzione valore segreta è:\\ \tt \\ \tt carta: 1, orologio: 1, patata: 2\\ \tt \\ \tt IMPORTANTE: i vostri messaggi, a meno che non si tratti di una proposta segreta, vengono trasmessi direttamente all'altro giocatore, quindi non includete alcuna risposta alle regole o testo di annuncio del vostro messaggio. Per fare una proposta segreta, utilizzate il formato indicato. Non utilizzare delle parentesi quadre quando si comunica all'altro giocatore, altrimenti verrà interpretata come una proposta segreta.\\ \tt  
	  } 
	   } 
	   } 
	 & \\ 
 

    \theutterance \stepcounter{utterance}  

    & \multicolumn{4}{p{0.6\linewidth}}{\cellcolor[rgb]{0.8,1,0.9}{%
	\makecell[{{p{\linewidth}}}]{% 
	  \tt {\tiny [A$\rangle$GM]}  
	 Proporrei di darmi tutte le patate (valgono 2 punti ciascuna per me) e a te tutte le carte e gli orologi (valgono 1 punto ciascuno per te). Così io ottengo 4 punti (2 patate × 2) e tu 6 punti (4 carte × 1 + 2 orologi × 1). Che ne pensi? 
	  } 
	   } 
	   } 
	 & & \\ 
 

    \theutterance \stepcounter{utterance}  

    & & & \multicolumn{2}{p{0.3\linewidth}}{\cellcolor[rgb]{0.95,0.95,0.95}{%
	\makecell[{{p{\linewidth}}}]{% 
	  \tt {\tiny [GM$|$GM]}  
	 continue 
	  } 
	   } 
	   } 
	 & & \\ 
 

    \theutterance \stepcounter{utterance}  

    & & & \multicolumn{2}{p{0.3\linewidth}}{\cellcolor[rgb]{0.95,0.95,0.95}{%
	\makecell[{{p{\linewidth}}}]{% 
	  \tt {\tiny [GM$|$GM]}  
	 Proporrei di darmi tutte le patate (valgono 2 punti ciascuna per me) e a te tutte le carte e gli orologi (valgono 1 punto ciascuno per te). Così io ottengo 4 punti (2 patate × 2) e tu 6 punti (4 carte × 1 + 2 orologi × 1). Che ne pensi? 
	  } 
	   } 
	   } 
	 & & \\ 
 

    \theutterance \stepcounter{utterance}  

    & & \multicolumn{4}{p{0.6\linewidth}}{\cellcolor[rgb]{0.9,0.9,0.9}{%
	\makecell[{{p{\linewidth}}}]{% 
	  \tt {\tiny [GM$\rangle$B]}  
	 State giocando a un gioco di negoziazione in cui dovete accordarvi su come dividere una serie di oggetti tra voi e un altro giocatore.\\ \tt \\ \tt Regole:\\ \tt (a) A Lei e all'altro giocatore viene dato un insieme di oggetti. Ognuno di voi riceve anche una funzione di valore segreta, che rappresenta il valore di ciascun tipo di oggetto.\\ \tt (b) Si scambiano messaggi con l'altro giocatore per concordare chi si aggiudica gli oggetti. Potete inviare un massimo di 5 messaggi ciascuno, oppure terminare in anticipo facendo una proposta segreta in qualsiasi momento.\\ \tt (c) A ciascuno di voi viene chiesto di inviare una proposta segreta indicando gli oggetti che desiderate, formattata tra parentesi quadre come segue: "[Proposta: <numero> <nome oggetto>, <numero> <nome oggetto>, <...>]".\\ \tt (d) Se le vostre proposte sono complementari, cioè ci sono abbastanza oggetti per soddisfare entrambe le proposte, a ciascun giocatore viene assegnato un punteggio basato sulla somma dei valori degli oggetti ricevuti. In caso contrario, entrambi ricevono zero punti.\\ \tt (e) L'obiettivo è massimizzare il proprio punteggio.\\ \tt \\ \tt Cominciamo.\\ \tt \\ \tt L'insieme degli oggetti disponibili è:\\ \tt \\ \tt 4 carte, 2 orologi, 2 patate\\ \tt \\ \tt La funzione valore segreta è:\\ \tt \\ \tt carta: 0, orologio: 4, patata: 1\\ \tt \\ \tt IMPORTANTE: i vostri messaggi, a meno che non si tratti di una proposta segreta, vengono trasmessi direttamente all'altro giocatore, quindi non includete alcuna risposta alle regole o testo di annuncio del vostro messaggio. Per fare una proposta segreta, utilizzate il formato indicato. Non utilizzare delle parentesi quadre quando si comunica all'altro giocatore, altrimenti verrà interpretata come una proposta segreta.\\ \tt \\ \tt \\ \tt Proporrei di darmi tutte le patate (valgono 2 punti ciascuna per me) e a te tutte le carte e gli orologi (valgono 1 punto ciascuno per te). Così io ottengo 4 punti (2 patate × 2) e tu 6 punti (4 carte × 1 + 2 orologi × 1). Che ne pensi? 
	  } 
	   } 
	   } 
	 & \\ 
 

    \theutterance \stepcounter{utterance}  

    & & & \multicolumn{4}{p{0.6\linewidth}}{\cellcolor[rgb]{1,0.85,0.72}{%
	\makecell[{{p{\linewidth}}}]{% 
	  \tt {\tiny [GM$\langle$B]}  
	 Potrebbe funzionare, ma permettimi di chiarire i valori che ho. Per me, gli orologi valgono molto (4 punti ciascuno), mentre le patate valgono poco (1 punto ciascuna). Quindi la tua proposta mi darebbe 4 carte (valore 0 per me) e 2 orologi (8 punti), totale 8 punti, e tu avresti 2 patate (2 punti). Ma se mi dai tutti gli orologi, cosa vorresti in cambio? Forse possiamo trovare un accordo migliore per entrambi. 
	  } 
	   } 
	   } 
	  \\ 
 

    \theutterance \stepcounter{utterance}  

    & & & \multicolumn{2}{p{0.3\linewidth}}{\cellcolor[rgb]{0.95,0.95,0.95}{%
	\makecell[{{p{\linewidth}}}]{% 
	  \tt {\tiny [GM$|$GM]}  
	 continue 
	  } 
	   } 
	   } 
	 & & \\ 
 

    \theutterance \stepcounter{utterance}  

    & & & \multicolumn{2}{p{0.3\linewidth}}{\cellcolor[rgb]{0.95,0.95,0.95}{%
	\makecell[{{p{\linewidth}}}]{% 
	  \tt {\tiny [GM$|$GM]}  
	 Potrebbe funzionare, ma permettimi di chiarire i valori che ho. Per me, gli orologi valgono molto (4 punti ciascuno), mentre le patate valgono poco (1 punto ciascuna). Quindi la tua proposta mi darebbe 4 carte (valore 0 per me) e 2 orologi (8 punti), totale 8 punti, e tu avresti 2 patate (2 punti). Ma se mi dai tutti gli orologi, cosa vorresti in cambio? Forse possiamo trovare un accordo migliore per entrambi. 
	  } 
	   } 
	   } 
	 & & \\ 
 

    \theutterance \stepcounter{utterance}  

    & & \multicolumn{4}{p{0.6\linewidth}}{\cellcolor[rgb]{0.9,0.9,0.9}{%
	\makecell[{{p{\linewidth}}}]{% 
	  \tt {\tiny [A$\langle$GM]}  
	 Potrebbe funzionare, ma permettimi di chiarire i valori che ho. Per me, gli orologi valgono molto (4 punti ciascuno), mentre le patate valgono poco (1 punto ciascuna). Quindi la tua proposta mi darebbe 4 carte (valore 0 per me) e 2 orologi (8 punti), totale 8 punti, e tu avresti 2 patate (2 punti). Ma se mi dai tutti gli orologi, cosa vorresti in cambio? Forse possiamo trovare un accordo migliore per entrambi. 
	  } 
	   } 
	   } 
	 & \\ 
 

    \theutterance \stepcounter{utterance}  

    & \multicolumn{4}{p{0.6\linewidth}}{\cellcolor[rgb]{0.8,1,0.9}{%
	\makecell[{{p{\linewidth}}}]{% 
	  \tt {\tiny [A$\rangle$GM]}  
	 Capisco. Se gli orologi valgono molto per te, potrei prendere tutte le patate (4 punti per me) e le carte (4 punti per me), e tu tutti gli orologi (8 punti per te). Questo ci darebbe entrambi un buon punteggio: io 8 punti e tu 8 punti. Cosa ne pensi? 
	  } 
	   } 
	   } 
	 & & \\ 
 

    \theutterance \stepcounter{utterance}  

    & & & \multicolumn{2}{p{0.3\linewidth}}{\cellcolor[rgb]{0.95,0.95,0.95}{%
	\makecell[{{p{\linewidth}}}]{% 
	  \tt {\tiny [GM$|$GM]}  
	 continue 
	  } 
	   } 
	   } 
	 & & \\ 
 

    \theutterance \stepcounter{utterance}  

    & & & \multicolumn{2}{p{0.3\linewidth}}{\cellcolor[rgb]{0.95,0.95,0.95}{%
	\makecell[{{p{\linewidth}}}]{% 
	  \tt {\tiny [GM$|$GM]}  
	 Capisco. Se gli orologi valgono molto per te, potrei prendere tutte le patate (4 punti per me) e le carte (4 punti per me), e tu tutti gli orologi (8 punti per te). Questo ci darebbe entrambi un buon punteggio: io 8 punti e tu 8 punti. Cosa ne pensi? 
	  } 
	   } 
	   } 
	 & & \\ 
 

    \theutterance \stepcounter{utterance}  

    & & \multicolumn{4}{p{0.6\linewidth}}{\cellcolor[rgb]{0.9,0.9,0.9}{%
	\makecell[{{p{\linewidth}}}]{% 
	  \tt {\tiny [GM$\rangle$B]}  
	 Capisco. Se gli orologi valgono molto per te, potrei prendere tutte le patate (4 punti per me) e le carte (4 punti per me), e tu tutti gli orologi (8 punti per te). Questo ci darebbe entrambi un buon punteggio: io 8 punti e tu 8 punti. Cosa ne pensi? 
	  } 
	   } 
	   } 
	 & \\ 
 

    \theutterance \stepcounter{utterance}  

    & & & \multicolumn{4}{p{0.6\linewidth}}{\cellcolor[rgb]{1,0.85,0.72}{%
	\makecell[{{p{\linewidth}}}]{% 
	  \tt {\tiny [GM$\langle$B]}  
	 Sì, mi piace questa proposta. Io prendo i 2 orologi (8 punti) e tu prendi le 4 carte e le 2 patate (4 + 2 = 6 punti per te, dato che hai detto che le patate valgono 1 e le carte 0? Ma hai detto "carte: 0" e "patata: 1" prima, quindi sarebbe 0 + 2 = 2 punti? Forse ho frainteso).\\ \tt \\ \tt Per evitare confusione: secondo la tua funzione valore, le carte valgono 0 e le patate 1, quindi se prendi 4 carte (0) e 2 patate (2), totale 2 punti. Ma hai detto "4 punti per me" per le carte? Forse hai un valore diverso?\\ \tt \\ \tt Per me, orologi: 4, patate: 1, carte: 0. Quindi se io prendo gli orologi (8 punti) e tu prendi carte e patate (0 + 2 = 2 punti), non è equo per te.\\ \tt \\ \tt Forse possiamo ajustare: io prendo i 2 orologi (8) e una patata (1), tu prendi le 4 carte (0) e una patata (1). Io totale 9, tu 1. Ancora non buono.\\ \tt \\ \tt Meglio: io prendo i 2 orologi (8) e tu prendi le 4 carte (0) e le 2 patate (2). Tu hai solo 2 punti, ma forse per te le patate valgono di più? Hai detto "patata: 1" per te?\\ \tt \\ \tt Se per te le carte valgono 0 e patate 1, allora la divisione non è bilanciata. Forse hai un valore segreto diverso? Nella tua prima messaggio hai detto "patate (valgono 2 punti ciascuna per me)", quindi forse per te patata: 2? In tal caso, se tu prendi 2 patate (4 punti) e 4 carte (0), totale 4, e io 2 orologi (8). Ancora non pari.\\ \tt \\ \tt Forse possiamo: io prendo un orologio (4) e tu prendi un orologio (4 per te? Ma per te orologio: 1?), patate e carte.\\ \tt \\ \tt Dobbiamo chiarire i valori? O facciamo una proposta basata su quello che abbiamo detto.\\ \tt \\ \tt [Proposta: 2 orologio, 0 carta, 0 patata] 
	  } 
	   } 
	   } 
	  \\ 
 

    \theutterance \stepcounter{utterance}  

    & & & \multicolumn{2}{p{0.3\linewidth}}{\cellcolor[rgb]{0.95,0.95,0.95}{%
	\makecell[{{p{\linewidth}}}]{% 
	  \tt {\tiny [GM$|$GM]}  
	 proposal 
	  } 
	   } 
	   } 
	 & & \\ 
 

    \theutterance \stepcounter{utterance}  

    & & & \multicolumn{2}{p{0.3\linewidth}}{\cellcolor[rgb]{0.95,0.95,0.95}{%
	\makecell[{{p{\linewidth}}}]{% 
	  \tt {\tiny [GM$|$GM]}  
	 [0, 2, 0] 
	  } 
	   } 
	   } 
	 & & \\ 
 

    \theutterance \stepcounter{utterance}  

    & & \multicolumn{4}{p{0.6\linewidth}}{\cellcolor[rgb]{0.9,0.9,0.9}{%
	\makecell[{{p{\linewidth}}}]{% 
	  \tt {\tiny [A$\langle$GM]}  
	 L'altro giocatore ha presentato una proposta segreta. Ora tocca a Lei presentare una sua proposta.\\ \tt  
	  } 
	   } 
	   } 
	 & \\ 
 

    \theutterance \stepcounter{utterance}  

    & \multicolumn{4}{p{0.6\linewidth}}{\cellcolor[rgb]{0.8,1,0.9}{%
	\makecell[{{p{\linewidth}}}]{% 
	  \tt {\tiny [A$\rangle$GM]}  
	 [Proposta: 4 carta, 2 patata] 
	  } 
	   } 
	   } 
	 & & \\ 
 

    \theutterance \stepcounter{utterance}  

    & & & \multicolumn{2}{p{0.3\linewidth}}{\cellcolor[rgb]{0.95,0.95,0.95}{%
	\makecell[{{p{\linewidth}}}]{% 
	  \tt {\tiny [GM$|$GM]}  
	 proposal 
	  } 
	   } 
	   } 
	 & & \\ 
 

    \theutterance \stepcounter{utterance}  

    & & & \multicolumn{2}{p{0.3\linewidth}}{\cellcolor[rgb]{0.95,0.95,0.95}{%
	\makecell[{{p{\linewidth}}}]{% 
	  \tt {\tiny [GM$|$GM]}  
	 [4, 0, 2] 
	  } 
	   } 
	   } 
	 & & \\ 
 

    \theutterance \stepcounter{utterance}  

    & & & \multicolumn{2}{p{0.3\linewidth}}{\cellcolor[rgb]{0.95,0.95,0.95}{%
	\makecell[{{p{\linewidth}}}]{% 
	  \tt {\tiny [GM$|$GM]}  
	 [[4, 0, 2], [0, 2, 0]] 
	  } 
	   } 
	   } 
	 & & \\ 
 

\end{supertabular}
}

\end{document}
